% =========================
% tables/summary_table.tex
% IJMI-required Summary Table (2-4 bullet points)
% Uses \caption* (unnumbered) so it does not consume the table counter
% =========================

\begin{table}[H]
\caption*{\textbf{Summary Table}}
\centering
\small
\renewcommand{\arraystretch}{1.55}
\setlength{\tabcolsep}{10pt}
\arrayrulecolor{headerblue}
\begin{tabular}{p{14cm}}
\toprule
\rowcolor{headerblue}
\textcolor{white}{\textbf{What this systematic review adds}} \\
\midrule
\textcolor{headerblue}{\textbullet}\ eHealth adoption is determined by alignment across four interacting domains—User Readiness, Technical Interoperability, Institutional Alignment, and Organizational Execution—not by any single domain in isolation. \\
\textcolor{headerblue}{\textbullet}\ Across 501 included studies and four national implementation cases, most studies presented determinants spanning three or more domains and none identified a single-domain explanation as sufficient; multi-domain convergence emerged as the decisive predictor of adoption maturity. \\
\textcolor{headerblue}{\textbullet}\ Strong institutional mandate and substantial technical investment are necessary but insufficient: the VA/DoD case demonstrates how realised gaps in user readiness, technical interoperability, and organizational execution constrain adoption value despite congressional mandate and sustained investment. \\
\textcolor{headerblue}{\textbullet}\ The U-T-I-O Model functions as both a descriptive framework and a diagnostic tool, providing a structured basis for assessing implementation readiness and identifying domain-level misalignment risks before deployment. \\
\bottomrule
\end{tabular}
\arrayrulecolor{black}
\end{table}
